% Symmetrie, Abbildungen und Koordinatensystem – bank/symmetrie-abbildungen/ (zusammenbau v0.4, Heft)
\einheitenkopf[t9]{Symmetrie, Abbildungen und Koordinatensystem}
\setcounter{aufgabe}{61}

% Anlauf (ohne Prüfkennung)
\begin{aufgabe}{Punkte und Figuren im Koordinatensystem – Anlauf}
\begin{teile}
\teil Du sollst $(5|1)$ eintragen. Wohin gehst du vom Ursprung zuerst? \\ \kreuz{zuerst nach rechts} \\ \kreuz{zuerst nach oben}
\teil Trage den Punkt $A(4|6)$ ein.

\begin{ksys}[xmin=0,xmax=8,ymin=0,ymax=8] \end{ksys}
\teil Trage die Punkte $B(-4|1)$ und $C(2|-3)$ ein.

\begin{ksys}[xmin=-6,xmax=6,ymin=-6,ymax=6] \end{ksys}
\end{teile}
\end{aufgabe}

\begin{aufgabe}{Punkte und Figuren im Koordinatensystem – Prüfungsaufgaben}
\begin{teile}
\teil Genau einer der vier Punkte liegt auf der Rechtsachse (x-Achse). Welcher? \hfill (P10 2020 OS) \\ \kreuz{$A(6|4)$} \\ \kreuz{$B(0|4)$} \\ \kreuz{$C(-6|-4)$} \\ \kreuz{$D(6|0)$}
\teil Genau einer der vier Punkte liegt auf der Rechtsachse (x-Achse). Welcher? \hfill (P10 2020 OS) \\ \kreuz{$E(7|1)$} \\ \kreuz{$F(0|1)$} \\ \kreuz{$G(-7|-1)$} \\ \kreuz{$H(7|0)$}
\end{teile}
\end{aufgabe}

% Anlauf (ohne Prüfkennung)
\begin{aufgabe}{Symmetrieachsen bestimmen – Anlauf}
\begin{teile}
\teil Passen beim Falten an der eingezeichneten Geraden beide Hälften genau aufeinander? \\ \kreuz{ja} \\ \kreuz{nein}

\drachen[achsen=senkrecht]{4}{3}{0.3}
\teil Zeichne die Symmetrieachse der Figur ein.

\drachen{4}{3}{0.3}
\teil Welche Figur hat keine Symmetrieachse? \\ \kreuz{Rechteck} \\ \kreuz{Parallelogramm} \\ \kreuz{Raute}
\end{teile}
\end{aufgabe}

\begin{aufgabe}{Symmetrieachsen bestimmen – Prüfungsaufgaben}
\begin{teile}
\teil Die Platte eines Gruppentischs hat die Form eines gleichschenkligen Trapezes. Wie viele Symmetrieachsen hat sie? Kreuze an. \hfill (P10 2022 OS) \\ \kreuz{0} \\ \kreuz{1} \\ \kreuz{2} \\ \kreuz{3} \\ \kreuz{4} \\ \kreuz{5}

\trapez{6}{3}{2}
\teil Ein Dachfenster hat die Form eines gleichschenkligen Trapezes. Wie viele Symmetrieachsen hat es? Kreuze an. \hfill (P10 2022 OS) \\ \kreuz{0} \\ \kreuz{1} \\ \kreuz{2} \\ \kreuz{3} \\ \kreuz{4} \\ \kreuz{5}

\trapez{4}{2}{3}
\teil Eine Bodenfliese ist quadratisch. Wie viele Symmetrieachsen hat sie? Kreuze an. \hfill (P10 2021 OS) \\ \kreuz{0} \\ \kreuz{1} \\ \kreuz{2} \\ \kreuz{3} \\ \kreuz{4} \\ \kreuz{5}

\viereck{(0,0)}{(3.5,0)}{(3.5,3.5)}{(0,3.5)}
\teil Ein Sitzkissen ist quadratisch. Wie viele Symmetrieachsen hat es? Kreuze an. \hfill (P10 2021 OS) \\ \kreuz{0} \\ \kreuz{1} \\ \kreuz{2} \\ \kreuz{3} \\ \kreuz{4} \\ \kreuz{5}

\viereck{(0,0)}{(2.5,0)}{(2.5,2.5)}{(0,2.5)}
\teil Ein Kinderdrache hat die Form des Drachenvierecks ABCD. Seine Leisten sind die Diagonalen: AC ist 40\,cm, BD ist 70\,cm lang, sie kreuzen sich rechtwinklig. Wie viele Symmetrieachsen hat das Drachenviereck? Kreuze an. \hfill (P10 2025 OS) \\ \kreuz{0} \\ \kreuz{1} \\ \kreuz{2} \\ \kreuz{3} \\ \kreuz{4}

\drachen[diagonalen]{4}{2.3}{0.3}
\teil Ein Wimpel hat die Form des Drachenvierecks ABCD mit den Diagonalen AC = 30\,cm und BD = 48\,cm, die sich rechtwinklig kreuzen. Wie viele Symmetrieachsen hat das Drachenviereck? Kreuze an. \hfill (P10 2025 OS) \\ \kreuz{0} \\ \kreuz{1} \\ \kreuz{2} \\ \kreuz{3} \\ \kreuz{4}

\drachen[diagonalen]{4.8}{3}{0.35}
\end{teile}
\end{aufgabe}

% Anlauf (ohne Prüfkennung)
\begin{aufgabe}{Spiegeln und Spiegelbild benennen – Anlauf}
\begin{teile}
\teil Punkt P liegt 3 Kästchen links von der Achse, sein Bildpunkt 3 Kästchen rechts, beide auf gleicher Höhe. Ist richtig gespiegelt? \\ \kreuz{ja} \\ \kreuz{nein}
\teil Spiegle den Punkt P an der Geraden a. Zähle die Kästchen bis zur Achse.

\begin{ksys}[xmin=0,xmax=9,ymin=0,ymax=8] \asymptote{x=4}{a} \punkt{2}{3}{P} \end{ksys}
\teil Spiegle das Dreieck ABC an der Geraden a.

\begin{ksys}[xmin=0,xmax=10,ymin=0,ymax=7] \asymptote{x=4}{a} \punkt{1}{1}{A} \punkt{3}{1}{B} \punkt{2}{4}{C} \end{ksys}
\end{teile}
\end{aufgabe}

\begin{aufgabe}{Spiegeln und Spiegelbild benennen – Prüfungsaufgaben}
\begin{teile}
\teil Nora schneidet ein Dreieck mit drei verschieden langen Seiten aus Papier aus. Sie legt es mit einer Seite an die Gerade a und spiegelt es an a. Welches Viereck bilden Dreieck und Spiegelbild zusammen? \hfill (P10 2018 OS)
\teil Für ein Fensterbild wird ein Dreieck mit den Seiten 3\,cm, 4\,cm und 6\,cm an der 6\,cm langen Seite gespiegelt. Welches Viereck entsteht aus Dreieck und Spiegelbild? \hfill (P10 2018 OS)
\end{teile}
\end{aufgabe}
